\section{Introduction}\label{sec:intro}

\tk{You can skip this section as it's just old text that I've pasted here to reuse at some point but I'll refer to the jaws example concretely later}

Conformal factuality (CF)~\cite{mohri2024language} splits the response of an LLM into subclaims, scores each subclaim, and removes the subclaims whose score falls below a threshold that is calibrated on labelled queries, so that, with high probability, every remaining subclaim is supported. Conformal-RAG (C-RAG)~\cite{feng2025response} instead scores each subclaim by its relevance to documents retrieved with the query, so that the selection of subclaims is grounded in retrieved evidence. Both methods may abstain by removing every subclaim, which satisfies the guarantee but does not answer the query.

However, one key limiation remains. The subclaims can only be scored on what was actually retrieved, and the documents are retrieved with the question before the answer is written, so that valid subclaims may still end up being removed due to a lack of evidence. This problem is of particular concern when the answer to a user's question is spread across multiple documents and linked via a \emph{bridge entity}, that is, an entity that is not referred to in the original question. To illustrate this, take the following example query: ``Where was the director of Jaws born?''. Retrieval with this query will likely surface documents such as the article on the film Jaws, and the LLM's own knowledge may correctly identify the director as Steven Spielberg. His biography, however, is unlikely to be retrieved, since nothing in the question points to it. Suppose that the LLM produces two example subclaims, $c^{\star}_1 =$ ``Jaws was directed by Steven Spielberg'' and $c^{\star}_2 =$ ``Steven Spielberg was born in Cincinnati, Ohio''. Subclaim $c^{\star}_1$ is similar to the Jaws article and would likely be retained. Subclaim $c^{\star}_2$ is true, and the LLM states it correctly from its own knowledge, but no retrieved document mentions Spielberg's birthplace, and so its score is low and it would be removed. The filtered answer would then contain $c^{\star}_1$ but not $c^{\star}_2$, which is the subclaim that answers the question.
 
To address multi-hop questions, Rehman and Yeh \cite{rehman2026conformal} extend the Conformal-RAG methodology to a multi-hop retrieval setting, in which they perform multiple rounds of retrieval. After the first round, an LLM rewrites the query from the question and the evidence found so far, in the hope that bridge entities will be identified and that the accumulated evidence will better cover the specifics of the question. The scoring and the filtering remain the same as in Conformal-RAG. Their main result is that the guarantee carries over to multi-hop RAG, as stricter levels consistently increase the fraction of answers in which every retained subclaim is supported. At the 95\% level, however, only 4.41\%--31.09\% of the generated subclaims are retained, and only 9.70\%--51.40\% of the answers are non-empty, which means that much of the reliability is obtained by abstaining. 
 
There are two reasons why multi-hop retrieval of this kind may not recover the missing evidence. Firstly, the rewritten queries are driven by the question, not by the subclaims, which means that the retrieval has no reason to look for the evidence of a particular subclaim that is failing. Secondly, the score of C-RAG multiplies by the similarity of each document to the original question. As a result, a document about the bridge entity, such as Spielberg's biography, is penalised even when it is retrieved, as it is not very similar to the question.

There also exist active RAG methods~\cite{jiang2023active,asai2024selfrag}, which retrieve during or after generation, although they revise sentences of the answer rather than subclaims directly. For example, FLARE \cite{jiang2023active} retrieves whenever the next sentence of the answer contains low-confidence tokens and regenerates that sentence, and RARR \cite{gao2023rarr} retrieves evidence for a finished answer and revises the parts that the evidence does not support. Reworking parts of an answer in this way has therefore been considered before, but neither method gives a statistical guarantee on the factuality of the returned answer, as CF and C-RAG do.

In this work, we integrate active RAG with conformal factuality, so that the system can retrieve further evidence for the subclaims that it cannot yet support before the answer is filtered, while the returned answer keeps the guarantee of CF and C-RAG. \tk{Contributions and outline to be added.}